% !TEX root = ../agr-paper.tex
% ---------------------------------------------------------------------
% AGR-led framing: the system is the contribution, the ablation is
% validation. ~950 words, compressed from the thesis introduction.
%
% The one thing this section must not do is promise that the
% verification layer raises accuracy. It does not -- the precision column
% does not move when the layer is removed (Sec 6). What it delivers is
% the output contract. Claiming that here, and only that, is what keeps
% Section 6 from reading as a retraction.
%
% And claim the contract at its real width. Only the traversed-adjacency
% route attaches evidence, and the logger keeps the COUNT and discards
% the list. Both bounds are in this paper already -- Sec 3 for the
% routes, Sec 7 for the record -- so the introduction must not say
% anything Sec 7 has to take back. In particular do not write that the
% reasoning is "auditable" without saying auditable WHERE; Sec 7's
% heading is that the contract cannot be audited from the run records.
%
% Register, as of the September 2026 pass: short sentences, active voice
% where it reads at least as well, few em-dashes. The thesis chapters
% were rewritten to this register first; this file follows them.
% ---------------------------------------------------------------------

\section{Introduction}\label{sec:intro}

Large language models answer factual questions fluently. Often enough,
they answer them wrongly. The errors are systematic rather than random,
and they are hard to catch from the output alone, because a fabricated
answer reads in the same register as a correct one. Huang et
al.~\cite{huang2025survey} treat this as a consequence of how the models
are built, not as an incidental defect to be patched away. The problem
sharpens once an answer requires composing several facts rather than
recalling one. Errors compound: a wrong resolution at the first hop
rarely announces itself downstream. And a correctness check on the final
answer never sees the intermediate steps. So a system that answers
correctly by accident looks the same as one that answers correctly by
sound reasoning.

Retrieval-augmented generation~\cite{lewis2020retrieval} addresses the
first half of this problem by grounding the model in retrieved text. It
does not address the second. Retrieval runs once, before reasoning
begins, so a multi-hop question must be answered from whatever that
single query returned. The passage holding the \emph{intermediate}
entity need not resemble the original question at all, and the question
is the only signal the retriever can rank on. Representing the corpus as
a knowledge graph~\cite{bollacker2008freebase} turns composition into a
traversal rather than an inference over prose.
GraphRAG~\cite{edge2024local} exploits this, but it still retrieves
statically. A neighbourhood fetched in one shot cannot know what the
next reasoning step will need. It is either too small to contain the
answer or so large that the relevant edges are buried.

Agentic systems close that gap by interleaving the two steps. The model
takes one step, observes what it found, and decides where to look
next~\cite{sun2024think,chen2024plan,luo2024reasoning}. Retrieval
becomes navigation, and the system becomes an agent operating on the
graph rather than a pipeline reading from it. \Cref{sec:related} traces
that line of work.

One property survives this progression unchanged. Every system named
above emits whatever its final generation call produces from the
traversed context. Nothing checks that what the answer \emph{asserts} is
supported by what the agent retrieved. The gap is not fabrication. Our
measurements confirm that graph-navigating systems assert entities that
exist in the graph and are reachable along the paths they walked
(\Cref{sec:groundedness}). The gap is precision. An entity can sit in
the retrieved subgraph, have been retrieved correctly, and still be the
wrong answer: an intermediate node from the reasoning chain, a container
where the question asked for a member, or one of several plausible
neighbours chosen arbitrarily. A system that verifies retrieval but
never verifies assertion cannot catch any of these.

This paper presents \textbf{Agentic Graph Reasoning} (AGR), a
knowledge-graph question-answering framework built as an explicit state
machine over a constrained, deterministic graph tool API. The program
orchestrates the loop. The language model is consulted only at bounded
decision points. AGR's distinguishing component is a \emph{Structural
Verification Layer}. Before an answer is emitted, the draft is
decomposed into atomic claims, and each claim is checked against the
triples the agent actually traversed. The word \emph{structural} bounds
that check as much as it names it. The two graph routes ask whether an
edge joins a claim's endpoints, not whether it is the edge the claim
asserts. So a claim that X is Y's mother is certified by an edge
recording that X is Y's child, and only the entailment fallback reads
the relation. \Cref{sec:limits-contribution} reports the exposure that
follows. Claims that cannot be grounded trigger targeted re-exploration
or are removed. The claims the traversal itself grounds carry their
triples back with the answer.

We are precise about what that layer buys, because the ablation of
\Cref{sec:attribution} is precise about what it does not. Claim
verification does not measurably change how often AGR is right. It
changes how often AGR is right \emph{for a reason it can exhibit}. The
contribution is the output contract, an answer together with the
traversed evidence for it, rather than an accuracy gain.
\Cref{sec:attribution} reports the null that establishes this at the
same length as the effects it does detect.

On $400$ certified questions from each of WebQSP~\cite{yih2016value} and
ComplexWebQuestions~\cite{talmor2018web}, over a Freebase-derived
environment of $2.59$ million entities and $8.31$ million triples, AGR
reaches Hits@1 of $0.755$ and $0.522$. Four baselines share one frozen
backbone and one call budget with it. AGR leads all four, at roughly
half the language-model calls of the strongest. On ComplexWebQuestions
it is the only system whose accuracy rises with hop count. Its advantage
over the agentic baseline is located rather than asserted. On the
questions where the shared cap lets that baseline finish, the baseline
is slightly ahead. The entire aggregate margin comes from the $29\%$ and
$44\%$ of questions on which the cap truncates it.

% elsarticle supplies the period after a \paragraph title; a period
% inside the braces printed "Contributions.." in every build until
% September 2026.
\paragraph{Contributions} This paper contributes:

\begin{enumerate}
  \item \textbf{AGR}, a state-machine agent over a constrained graph tool
    API whose answers carry the traversed triples that support them, so
    the reasoning is auditable at emission rather than merely asserted
    (\Cref{sec:framework}). \Cref{sec:limits-contribution} bounds how
    far that reaches.
  \item A \textbf{Structural Verification Layer} that checks atomic
    claims against traversed adjacency before emission. Its repair
    routes and their firing rates are reported as they behave rather
    than as designed (\Cref{sec:verification}).
  \item A \textbf{controlled comparison} against four baselines on one
    frozen backbone under one call budget, so that differences are
    attributable to architecture rather than to model capacity
    (\Cref{sec:results}). Beside it we report a direct comparison
    against the published figures of the system whose data distribution
    we share, and that comparison does not favour this work
    (\Cref{sec:rog}).
  \item A \textbf{component-level ablation} with paired significance
    testing. It finds that removing the planner \emph{improves} accuracy
    on the shallower benchmark while reducing token use, and that three
    of four components show no detectable effect
    (\Cref{sec:attribution}).
  \item A \textbf{hand-read census} of all $259$ remaining failures. It
    names the \emph{echo attractor}, a true, grounded entity one hop
    from the answer that every grounding check passes, and it
    quantifies label defects in both benchmarks
    (\Cref{sec:erroranalysis}).
\end{enumerate}

\Cref{sec:related} positions AGR against prior work.
\Cref{sec:framework} specifies the framework and the verification layer.
\Cref{sec:setup} describes the environment and the experimental
protocol, and \Cref{sec:results} reports the measurements.
\Cref{sec:attribution} attributes them to components,
\Cref{sec:erroranalysis} reads what remains, and \Cref{sec:discussion}
states the limits.
